%% SUPERSEDED by p3-intro-related-methodology.tex (the P3 deliverable), which
%% carries these sections forward in reconciled form and adds Section III.
%% Kept as the P2 snapshot; do not edit further.
%% Trust-Gated Fast-Weight Updates for TTT-E2E LLMs
%% Introduction, Related Work, and References
%% Self-contained: compiles with a stock IEEEtran installation (e.g. Overleaf),
%% no external .bib file required.

\documentclass[journal]{IEEEtran}

\usepackage{cite}
\usepackage{amsmath}
\usepackage{array}

\begin{document}

\title{Trust-Gated Fast-Weight Updates for TTT-E2E Large Language Models}

\author{Manas~Maahir~R.,~Jay~Krishna~Kamlekar,~and~Kartheek~Yadav% <-this % stops a space
\thanks{The authors are with the School of Computer Science and Engineering,
Vellore Institute of Technology, Vellore, India.}}

\maketitle

\section{Introduction}

\IEEEPARstart{L}{arge} language models are increasingly deployed on inputs far
longer than the contexts they were trained to attend over, and the quadratic cost
of full attention makes naive scaling impractical. A recent line of work responds
by treating long-context modeling not as an architecture problem but as a learning
problem: rather than attending over a context, the model learns it. End-to-end
test-time training (TTT-E2E) instantiates this idea directly~\cite{tandon2025tttE2E}.
The model carries a set of \emph{fast weights} that it updates by next-token
prediction as it reads, compressing the context into parameters at near-recurrent
cost while matching the scaling behavior of full attention. Adaptation at inference
is no longer a preprocessing step or an auxiliary objective; it is the mechanism by
which the model represents what it has read. The same principle now appears across
deployed systems that personalize, accumulate session state, or maintain persistent
memory, and its practical appeal is clear: capability that would otherwise demand a
longer context window or a larger model is obtained instead by letting the model
keep learning while it serves.

That capability carries a security cost which has not been accounted for. Once fast
weights are updated during inference, the serving path becomes a \emph{write path}
into model state, and a system that learns while it serves can be poisoned while it
serves. The concern is not hypothetical. Test-time adaptation has been shown to be
poisonable under white-box assumptions~\cite{cong2024ttp}, under realistic
constraints on what an attacker can actually submit~\cite{su2025realistic}, and in
the black-box continual setting~\cite{hoang2024rip}. What makes the fast-weight case
difficult is the absence of nearly every signal on which conventional defenses rely.
Updates are self-supervised, so there is no label against which to check them. They
arrive from a single stream, so there is no cohort of peer updates whose agreement
could expose an outlier. Each individual update is small and, taken alone, entirely
plausible; the damage is cumulative, distributed across a sequence of steps that a
per-update test would pass one at a time. And because the appeal of the approach is
its near-recurrent cost, a defense that inspects every update must be cheap enough
not to erase the advantage it is protecting.

Existing defenses were built for settings that differ from this one structurally,
not merely superficially. Training-time methods harden a model before deployment,
whether by adversarial self-training with consistency regularization across weakly
augmented and adversarial views~\cite{li2024ast}, by routing training points
according to estimated poisoning risk~\cite{bena2025risk}, or by distilling a clean
student from a suspect encoder~\cite{han2025mimic}; all act on slow weights that are
frozen once serving begins. Defenses for distributed learning do inspect updates
before they commit, using robust aggregation over client
gradients~\cite{yang2025scalesign,zukaib2026d3flids}, an audit against a clean
reference dataset~\cite{basak2025dpad,nawshin2026concurrent}, or a learned selection
policy~\cite{xu2026feddrlpd} --- but each depends on a population of concurrent
updates whose mutual agreement carries the signal. Robust test-time adaptation
offers the closest analogue, replacing the mean in batch-normalization statistics
with the median so that no single malicious sample can swing
them~\cite{park2024medbn}. Finally, work on contextual backdoors establishes that
adaptation-time input is already an exploitable surface~\cite{liu2025contextual},
though it changes no parameters at all.

The resulting gap is specific rather than general. No published mechanism governs
\emph{parametric} updates committed at inference time in a language model. The
closest defense, MedBN~\cite{park2024medbn}, operates on batch-normalization
statistics, and the TTT-E2E architecture contains no batch-normalization layers for
such a test to act upon. Robust aggregation presupposes the redundancy of many
concurrent clients, which a single context stream does not provide. Audit-based
gates presuppose both a clean reference corpus and a discrete round boundary at
which to run the audit, and neither exists when updates commit token by token. The
threat itself has been established in image-model test-time
adaptation~\cite{cong2024ttp,su2025realistic,hoang2024rip} and has not been carried
across to fast-weight language models, where the unit under attack is neither a
normalization statistic nor an aggregated gradient but the weight delta itself.
Just as consequentially, no existing method bounds how far accumulated updates may
carry a model from its trusted state, and none provides a route back once they have.

This paper addresses that gap. We propose a \emph{trust gate} for fast-weight
updates: a mechanism that intercepts every proposed update at the point where it
would commit, admits it only if it remains consistent with a frozen anchor model,
holds the accumulated effect of admitted updates within an explicit drift budget
$\varepsilon$, and restores the last trusted checkpoint when that budget is
breached. The objective is a defended inference path whose divergence from trusted
behavior is bounded by construction rather than observed after the fact.

The main contributions of this paper are:

\begin{itemize}
\item We formalize benign-stream fast-weight poisoning as a threat model for
TTT-E2E language models, and construct a crafted-stream adversary that steers fast
weights using inputs which remain individually unremarkable, giving the defense a
concrete attacker to be evaluated against.

\item We design a trust gate that intercepts each proposed fast-weight update at
its single commit point and scores it on two signals: behavioral consistency with a
frozen meta-learned anchor over a rotating held-out probe set, and the uncertainty
of the proposed update itself.

\item We introduce a latching cumulative-drift accumulator that converts per-update
decisions into a system-level property, and we establish a bounded per-window drift
guarantee: over any window, the sequence of admitted updates cannot displace the
fast weights beyond the configured budget.

\item We show that the enforcement primitives are $O(1)$ in the streaming setting.
Because fast-weight updates are pure functions over a parameter tree, rejection
reduces to an elementwise select between weight trees and rollback to a reference
swap into a ring buffer of checkpoints, with no mutable state to unwind.

\item We specify an evaluation protocol in which corruption thresholds are fixed
before any measurement is taken, comparing the gate against a robust-aggregation
baseline constructed as the closest available analogue of MedBN for this setting,
and recording every decision in an accept/reject/rollback audit trail.
\end{itemize}

The significance of this work is both theoretical and practical. Theoretically, it
supplies a drift guarantee for inference-time parametric adaptation, a setting in
which robustness has so far been argued empirically rather than bounded.
Practically, the gate is an \emph{overlay}: it wraps the update function rather than
modifying the model, so it requires neither retraining nor any change to published
weights, and can be applied to an existing deployment. More broadly, any system that
writes to its own parameters while serving --- a personalized assistant, an agent
with persistent memory, a continually adapting recommender --- inherits the same
write-path exposure, and the decomposition into interceptor, drift budget, and
rollback is not specific to TTT-E2E.

The remainder of this paper is organized as follows. Section~II reviews existing
work on poisoning and defense across training-time, distributed, structural,
context-level, and test-time adaptation settings, and identifies the gap this work
addresses. Section~III presents the threat model and the trust-gate architecture.
Section~IV describes the experimental setup and the pre-registered evaluation
criteria. Section~V reports results, and Section~VI concludes.

\section{Related Work}
\label{sec:related}

\subsection{Thematic Organization of Existing Work}

Research relevant to the poisoning of adaptive models divides into five families,
distinguished less by application domain than by \emph{where in a model's lifecycle
an adversary writes} and \emph{what unit a defense inspects}. These are: (i)
poisoning of training data and offline hardening against it; (ii) update-level
defense in distributed learning, where a server screens proposed parameter changes
before aggregating them; (iii) structural and feature poisoning in relational
models, where individually plausible perturbations accumulate; (iv) inference-time
context poisoning, which alters model behavior without touching parameters; and (v)
attacks on, and defenses for, test-time adaptation --- the family closest to the
setting considered here. Read in order, these families trace a steady migration of
the attack surface, from the training corpus to the update channel to the input
stream at inference. The defenses have not migrated at the same rate.

\subsection{Key Representative Studies}

In the training-time family, Bena \emph{et al.} score each training point with three
geometric indicators of poisoning and route high-risk points across an ensemble
rather than deleting them, avoiding the false-positive cost of
filtering~\cite{bena2025risk}. Han \emph{et al.} address the supply-chain case,
distilling a randomly initialized student from a suspect self-supervised encoder and
using per-layer mutual information to decide which layers carry clean signal worth
copying~\cite{han2025mimic}. Li \emph{et al.} act on the model rather than the data,
using confident pseudo-labels from weakly augmented views as targets for adversarial
views so that representations remain consistent under perturbation~\cite{li2024ast}.

The distributed family is the largest and the most directly concerned with screening
updates. Yang \emph{et al.} contribute both sides of the exchange: ScaleSign builds
a malicious gradient by per-coordinate scaling that preserves norm, direction, and
sign statistics simultaneously, and MSGuard answers it by clustering clients over
five features rather than one~\cite{yang2025scalesign}. Zukaib \emph{et al.}
dispense with the auxiliary dataset entirely, normalizing gradient magnitudes across
layers and projecting misaligned directions onto a consensus subspace before an
isolation forest screens what remains~\cite{zukaib2026d3flids}. A second group gates
on behavior rather than geometry: Basak and Chatterjee evaluate each update against
a public audit set and discard those exceeding a divergence
tolerance~\cite{basak2025dpad}, while Nawshin \emph{et al.} clamp update norms and
then require that a client's deviation in behavioral space agree with its deviation
in parameter space~\cite{nawshin2026concurrent}. Xu \emph{et al.} replace fixed
thresholds with a learned policy, training an agent to select trustworthy clients
from observed system state~\cite{xu2026feddrlpd}. Liu \emph{et al.} push detection
into the encrypted domain, amplifying benign--malicious separation under homomorphic
encryption~\cite{liu2025antidotefl}. Huang \emph{et al.} sidestep identification
altogether, estimating the \emph{fraction} of contaminated reports from cross-round
consistency and subtracting its contribution from the
aggregate~\cite{huang2024ldpguard}.

Within relational models, Shi \emph{et al.} inject statistically plausible nodes
into hypergraphs under an explicit cost budget, searching hyperedge combinations
with a genetic algorithm~\cite{shi2025h3ni}, while Chen \emph{et al.} accumulate
feature gradients with momentum so that a sequence of individually budgeted
modifications compounds into large degradation~\cite{chen2025mghga}. On the defense
side, Zhang and Bao enforce consistency between a purified local view and a
diffusion-based global view of the same graph~\cite{zhang2025cola}.

The context-poisoning family is represented by Liu \emph{et al.}, who show that
poisoning a handful of in-context demonstrations induces a closed-box language model
to emit code containing a latent defect, activated only when a textual trigger
appears in the prompt and a matching visual trigger appears in the agent's
environment~\cite{liu2025contextual}. No weights change at any point.

Finally, in test-time adaptation, Cong \emph{et al.} established that adaptation is
poisonable~\cite{cong2024ttp}, Su \emph{et al.} constrained the attacker to
perturbations that could realistically be submitted~\cite{su2025realistic}, and
Hoang \emph{et al.} extended the attack to the black-box continual
setting~\cite{hoang2024rip}. Against these, Park \emph{et al.} proposed MedBN,
substituting the median for the mean in batch-normalization statistics so that a
single malicious sample in a test batch cannot swing them~\cite{park2024medbn}.
Wang \emph{et al.}, though targeting oriented object detection rather than
adaptation, contribute a result that bears on defense design: their attack ablation
shows individually modest search enhancements composing into a substantially
stronger attack~\cite{wang2026insulator}.

\subsection{Strengths and Limitations}

Six conclusions follow from this body of work, each bearing directly on the design
of a defense for fast weights.

First, a defense that inspects one signal is defeated by an attacker who matches
that signal. ScaleSign is the cleanest demonstration: by preserving magnitude,
direction, and sign statistics at once, it is accepted in the overwhelming majority
of rounds by a defense that examines exactly those
quantities~\cite{yang2025scalesign}. The field's response has been to add features
--- a spectral score, a projection step, a clustering stage~\cite{yang2025scalesign,
zukaib2026d3flids} --- which raises the attacker's cost without closing the loop.
These multi-feature defenses are moreover typically evaluated against attacks that
either predate them or originate with the same authors, so the question of how many
independent signals are needed before evasion becomes genuinely expensive remains
open.

Second, behavioral gating works, but imports an assumption that often goes
unexamined. Both the audit-based~\cite{basak2025dpad} and
cross-space~\cite{nawshin2026concurrent} defenses require the defender to hold a
clean reference set representative of honest behavior. Where legitimate data is
genuinely heterogeneous, such a set may reject valid updates as readily as malicious
ones; where data is private, it may not be obtainable at all. The design principle
survives the criticism --- comparing a proposed update against trusted reference
behavior is the right instinct, and it recurs in the consistency objectives
of~\cite{li2024ast} and~\cite{zhang2025cola} --- but the reference must be something
the defender actually possesses.

Third, per-sample detection has a hard limit. Huang \emph{et al.} make it concrete:
when a poisoned contribution passes through the same generative pipeline as a clean
one, no per-report test can separate them, and every prior countermeasure in their
setting is simply inapplicable~\cite{huang2024ldpguard}. Their answer --- estimate
contamination at the population level and correct the aggregate rather than
identifying culprits --- is a template worth naming, although it assumes a static
adversary and knowledge of the target set. The lesson generalizes beyond its
setting: where individual contributions are indistinguishable by construction, a
defense must reason about aggregate effect rather than individual identity.

Fourth, accumulated small perturbations are the operative threat. Momentum-guided
feature poisoning~\cite{chen2025mghga} and budget-aware node
injection~\cite{shi2025h3ni} both establish that a sequence of individually
plausible, individually budgeted changes compounds into large behavioral change, and
the attack ablation of~\cite{wang2026insulator} makes the same point from the
opposite direction. A defense that evaluates each update in isolation, and passes
each one, is not thereby safe: the quantity requiring constraint is cumulative
displacement, not per-step magnitude.

Fifth, defenses transfer across settings far less readily than their framing
suggests. Han \emph{et al.} quantify this directly, showing that defenses designed
for supervised classifiers and ported to self-supervised encoders leave a wide band
of residual attack success~\cite{han2025mimic}. The field was not merely
under-explored in the new setting; it was actively mis-served by transplanted
methods. This is a standing caution against assuming that a robust test-time
adaptation defense can simply be repointed at a new architecture.

Sixth, the defenses closest in spirit are all offline. Consistency regularization
against a trusted view~\cite{li2024ast}, multi-view
agreement~\cite{zhang2025cola}, and reconstruction-based
denoising~\cite{wang2026insulator} each reduce adversarial harm, but each operates
on slow weights during training or on inputs prior to inference. None intercepts a
parameter update as it commits, none bounds cumulative displacement from a trusted
state, and none provides a recovery path once displacement has occurred.

\subsection{The Research Gap}

Taken together, these limitations describe a gap that is structural rather than
incremental: \emph{no existing defense governs parametric updates committed at
inference time in a language model, and none bounds the cumulative drift such
updates induce.} The distinction from MedBN~\cite{park2024medbn} deserves precise
statement, since it is the nearest published defense and the one a reader is most
likely to raise. Table~\ref{tab:gap} sets out the five axes on which the settings
diverge.

\begin{table}[!t]
\renewcommand{\arraystretch}{1.3}
\caption{Robust Test-Time Adaptation Versus the Proposed Trust Gate}
\label{tab:gap}
\centering
\footnotesize
\begin{tabular}{|p{0.20\columnwidth}|p{0.32\columnwidth}|p{0.32\columnwidth}|}
\hline
\textbf{Axis} & \textbf{Robust TTA (MedBN)} & \textbf{Proposed trust gate} \\
\hline
Domain & Test-time adaptation for image classification & Long-context language
modeling with fast weights \\
\hline
Unit defended & Batch-normalization statistics & The fast-weight update
$\Delta\theta$ itself \\
\hline
Signal type & Distributional: a robust statistic over a test batch & Behavioral:
output divergence from a frozen anchor \\
\hline
Guarantee & Robust statistic within one batch & Bounded cumulative drift per window
\\
\hline
Recovery & None & Ring-buffered checkpoints with $O(1)$ rollback \\
\hline
\end{tabular}
\end{table}

The differences are not cosmetic. MedBN robustifies batch-normalization statistics,
and the TTT-E2E architecture contains no batch-normalization layers at all, so there
is no statistic on which a median-substitution test could operate. A robust statistic
bounds the influence of one sample within one batch; it says nothing about
displacement accumulated across a stream, which is precisely the regime the
cumulative-perturbation attacks~\cite{shi2025h3ni,chen2025mghga} identify as
dangerous. And MedBN, like the robust-aggregation family generally, has no notion of
checkpointing or rollback: a defense that admits an update it should have rejected
has no route back.

The consequences of leaving this gap open scale with adoption. The property that
makes adaptation at inference attractive --- that the model absorbs its input into
its parameters --- is exactly the property that makes the input stream a write
channel. As long as no mechanism governs that channel, every deployment of a
fast-weight model accepts unbounded, unlogged, and irreversible modification of its
own state by whoever supplies its context.

\subsection{Positioning of This Work}

Motivated by these limitations, this paper proposes a trust gate for fast-weight
updates in TTT-E2E language models. From the behavioral-gating
line~\cite{basak2025dpad,nawshin2026concurrent} we take the principle that an update
should be validated against trusted reference behavior before it commits, and we
resolve the clean-reference-set assumption by anchoring against a frozen copy of the
model itself, evaluated on rotating held-out probes, rather than against a corpus the
defender must independently obtain. From the cumulative-perturbation
attacks~\cite{shi2025h3ni,chen2025mghga,wang2026insulator} we take the requirement
that the constrained quantity be aggregate displacement, and we make it explicit as a
latching drift budget rather than leaving it as an emergent property. From the
aggregate-correction template of~\cite{huang2024ldpguard} we take the insight that
where individual contributions cannot be separated, a defense must reason about
accumulated effect. And from the transfer failures documented
in~\cite{han2025mimic} we take the discipline of not assuming that an existing
defense can be repointed at a new architecture: the gate is designed with the
fast-weight update as its unit, rather than adapted from a statistic that does not
exist in this setting. What no prior work supplies, and what this paper contributes,
is the combination of an interception point at the update commit, a bounded
cumulative-drift budget, and constant-time rollback to the last trusted checkpoint.

\begin{thebibliography}{20}

\bibitem{tandon2025tttE2E}
A. Tandon, K. Dalal, X. Li, D. Koceja, M. R\o{}d, S. Buchanan, X. Wang,
J. Leskovec, S. Koyejo, T. Hashimoto, C. Guestrin, J. McCaleb, Y. Choi, and Y. Sun,
``End-to-end test-time training for long context,'' \emph{arXiv:2512.23675}, 2025.

\bibitem{cong2024ttp}
T. Cong, X. He, Y. Shen, and Y. Zhang, ``Test-time poisoning attacks against
test-time adaptation models,'' in \emph{Proc. 45th IEEE Symp. Security and Privacy
(S\&P)}, 2024. % TODO: add page range from the published proceedings

\bibitem{su2025realistic}
Y. Su, Y. Li, N. Liu, K. Jia, X. Yang, C.-S. Foo, and X. Xu, ``On the adversarial
risk of test time adaptation: An investigation into realistic test-time data
poisoning,'' in \emph{Proc. Int. Conf. Learning Representations (ICLR)}, 2025.

\bibitem{hoang2024rip}
T.-H. Hoang, D. M. Vo, and M. N. Do, ``R.I.P.: A simple black-box attack on
continual test-time adaptation,'' \emph{arXiv:2412.01154}, 2024.

\bibitem{li2024ast}
Z. Li, M. Wu, C. Jin, D. Yu, and H. Yu, ``Adversarial self-training for robustness
and generalization,'' \emph{Pattern Recognition Letters}, vol. 185, pp. 117--123,
2024.

\bibitem{bena2025risk}
N. Bena, M. Anisetti, E. Damiani, C. Y. Yeun, and C. A. Ardagna, ``Protecting
machine learning from poisoning attacks: A risk-based approach,'' \emph{Computers
\& Security}, vol. 155, art. 104468, 2025.

\bibitem{han2025mimic}
T. Han, W. Sun, Z. Ding, C. Fang, H. Qian, J. Li, Z. Chen, and X. Zhang, ``Mutual
information guided backdoor mitigation for pre-trained encoders,'' \emph{IEEE Trans.
Information Forensics and Security}, vol. 20, pp. 3414--3428, 2025.

\bibitem{yang2025scalesign}
L. Yang, Y. Miao, Z. Liu, Z. Liu, X. Li, D. Kuang, H. Li, and R. H. Deng,
``Enhanced model poisoning attack and multi-strategy defense in federated
learning,'' \emph{IEEE Trans. Information Forensics and Security}, vol. 20,
pp. 3877--3892, 2025.

\bibitem{zukaib2026d3flids}
U. Zukaib, X. Cui, F. Niaz, D. Liang, C. Zheng, and S. U. Din, ``Defending federated
learning-based intrusion detection systems against model poisoning attacks,''
\emph{Neurocomputing}, vol. 698, art. 134208, 2026.

\bibitem{basak2025dpad}
S. Basak and K. Chatterjee, ``DPAD: Data poisoning attack defense mechanism for
federated learning-based system,'' \emph{Computers and Electrical Engineering},
vol. 121, art. 109893, 2025.

\bibitem{nawshin2026concurrent}
F. Nawshin, D. Unal, and P. N. Suganthan, ``Novel defense strategies for concurrent
data and model poisoning attacks in federated learning,'' \emph{Knowledge-Based
Systems}, vol. 345, art. 116057, 2026.

\bibitem{xu2026feddrlpd}
N. Xu, Y. Feng, N. Liu, M. Liu, Y. Li, and X. Fu, ``FedDRLPD: Deep reinforcement
learning-based defense mechanism against poisoning attacks in federated learning,''
\emph{Knowledge-Based Systems}, vol. 339, art. 115558, 2026.

\bibitem{park2024medbn}
H. Park, J. Hwang, S. Mun, S. Park, and J. Ok, ``MedBN: Robust test-time adaptation
against malicious test samples,'' in \emph{Proc. IEEE/CVF Conf. Computer Vision and
Pattern Recognition (CVPR)}, 2024. % TODO: add page range from the published proceedings

\bibitem{liu2025contextual}
A. Liu, Y. Zhou, X. Liu, T. Zhang, S. Liang, J. Wang, Y. Pu, T. Li, J. Zhang,
W. Zhou, Q. Guo, and D. Tao, ``Compromising LLM driven embodied agents with
contextual backdoor attacks,'' \emph{IEEE Trans. Information Forensics and
Security}, vol. 20, pp. 3979--3994, 2025.

\bibitem{liu2025antidotefl}
Y. Liu, H. Zhang, M. Wang, Q. Xie, and Z. Sun, ``AntidoteFL: Enhancing defense
against poisoning attacks in federated learning,'' \emph{Computer Networks},
vol. 269, art. 111427, 2025.

\bibitem{huang2024ldpguard}
K. Huang, G. Ouyang, Q. Ye, H. Hu, B. Zheng, X. Zhao, R. Zhang, and X. Zhou,
``LDPGuard: Defenses against data poisoning attacks to local differential privacy
protocols,'' \emph{IEEE Trans. Knowledge and Data Engineering}, vol. 36, no. 7,
pp. 3195--3209, 2024.

\bibitem{shi2025h3ni}
H. Shi, B. Zeng, R. Yu, Y. Yang, Z. Zouxia, C. Hu, and R. Shi, ``H3NI:
Non-target-specific node injection attacks on hypergraph neural networks via genetic
algorithm,'' \emph{Neurocomputing}, vol. 613, art. 128746, 2025.

\bibitem{chen2025mghga}
Y. Chen, S. Picek, Z. Ye, Z. Wang, and H. Zhao, ``Momentum gradient-based untargeted
poisoning attack on hypergraph neural networks,'' \emph{Neurocomputing}, vol. 634,
art. 129835, 2025.

\bibitem{zhang2025cola}
X. Zhang and P. Bao, ``Multi-view contrastive learning for graph adversarial
defense,'' \emph{Neural Networks}, vol. 192, art. 107868, 2025.

\bibitem{wang2026insulator}
Q. Wang, C. Yang, J. Du, N. Li, J. Wang, H. Han, Y. Hu, W. Peng, and J. Sun,
``Adversarial attack-defense framework for enhancing the robustness of power
insulator detection in cloud-edge deployment,'' \emph{Engineering Applications of
Artificial Intelligence}, vol. 166, art. 113625, 2026.

\end{thebibliography}

\end{document}
